\subsection{\texorpdfstring{\(\mathsf{T}(\mathbf{M})\) 与 \(\mathsf{T}(\mathbf{M}^{*})\)、\(\mathsf{S}(\mathbf{M})\) 与 \(\mathsf{S}(\mathbf{M}^{*})\)、\(\bigwedge(\mathbf{M})\) 与 \(\bigwedge(\mathbf{M}^{*})\) 之间的内乘}{T(M) 与 T(M*)、S(M) 与 S(M*)、wedge(M) 与 wedge(M*) 之间的内乘}}

右内乘在 \(\mathsf{T}(\mathbf{M})\) （相应地 \(\mathsf{S}(\mathbf{M})\) ，相应地 \(\bigwedge (\mathbf{M})\) ）上定义了代数 \(\mathsf{T}(\mathbf{M})^{*\mathfrak{g}\mathfrak{r}}\) （相应地 \(\mathsf{S}(\mathbf{M})^{*\mathfrak{g}\mathfrak{r}}\) ，相应地 \(\bigwedge (\mathbf{M})^{*\mathfrak{g}\mathfrak{r}}\) ）上的一个右模结构（第 7 号，例子）。利用第 5 号的典范同态 \(\theta_{\mathsf{T}}\) （相应地 \(\theta_{\mathsf{S}}\) ，相应地 \(\theta_{\wedge}\) ），我们得到

一个右 \(\mathsf{T}(\mathbf{M}^{*})\) -模结构，定义在 \(\mathsf{T}(\mathbf{M})\) 上

一个右 \(\mathbf{S}(\mathbf{M}^{*})\) -模结构，定义在 \(\mathbf{S}(\mathbf{M})\) 上

一个左 \(\bigwedge (\mathbf{M}^{*})\) -模结构，定义在 \(\bigwedge (\mathbf{M})\) 上。

这些结构中任何一个的外合成律也记为

\[
(z ^ {*}, t) \mapsto i (z ^ {*}). t
\]

（通过滥用语言）；我们也写作 \(t \llcorner z^*\) 以代替 \(i(z^*)\) . \(t\) ，在 \(\mathsf{T}(\mathbf{M})\) 或 \(\mathsf{S}(\mathbf{M})\) 的情形；另一方面，我们写作 \(z^* \lrcorner t\) ，在 \(\bigwedge (\mathbf{M})\) 的情形，并称其为 \(t\) 被 \(z^*\) 所作的左内乘，因为此时我们有一个左 \(\bigwedge (\mathbf{M}^*)\) -模法则。对齐次元素 \(z^*\) （次数为 \(n\) ）和齐次元素 \(t\) （次数为 \(p\) ），有 \(i(z^*)\) . \(t = 0\) （若 \(p < n\) ）；而对 \(x_i \in \mathbf{M} (1 \leqslant i \leqslant p)\) , \(x_j^* \in \mathbf{M}^* (1 \leqslant j \leqslant n)\) 及 \(p \geqslant n\) ，根据第 7 号的公式 (58)、(59) 和 (60)，

\[
\begin{array}{r l} i (x _ {1} ^ {*} \otimes x _ {2} ^ {*} \otimes \dots \otimes x _ {n} ^ {*}) \cdot (x _ {1} \otimes x _ {2} \otimes \dots \otimes x _ {p}) \\ & = \left(\prod_ {j = 1} ^ {n} \langle x _ {j} ^ {*}, x _ {j} \rangle\right) x _ {n + 1} \otimes \dots \otimes x _ {p} \end{array}\tag{66}
\]

\[
i \left(x _ {1} ^ {*} x _ {2} ^ {*} \dots x _ {n} ^ {*}\right). (x _ {1} x _ {2} \dots x _ {p}) = \sum_ {\sigma} \left(\prod_ {j = 1} ^ {n} \langle x _ {j} ^ {*}, x _ {\sigma (j)} \rangle\right) x _ {\sigma (n + 1)} \dots x _ {\sigma (p)}\tag{67}
\]

\[
\begin{array}{l} i \left(x _ {1} ^ {*} \wedge x _ {2} ^ {*} \wedge \dots \wedge x _ {n} ^ {*}\right). \left(x _ {1} \wedge x _ {2} \wedge \dots \wedge x _ {p}\right) \\ = (- 1) ^ {n (n - 1) / 2} \sum_ {\sigma} \varepsilon_ {\sigma} \left(\prod_ {j = 1} ^ {n} \langle x _ {j} ^ {*}, x _ {\sigma (j)} \rangle\right) x _ {\sigma (n + 1)} \wedge \dots \wedge x _ {\sigma (p)} \end{array}\tag{68}
\]

其中，在公式 (67) 和 (68) 中， \(\sigma\) 遍历排列 \(\sigma \in \mathfrak{S}_p\) 中在区间 \([1, n]\) 和 \([n + 1, p]\) 上递增的那些排列。我们也可以在内乘记号下写成

\[
\begin{array}{l l} \langle t \llcorner u ^ {*}, v ^ {*} \rangle = \langle t, \theta_ {\mathsf {T}} (u ^ {*} v ^ {*}) \rangle & \text {for} \quad t \in \mathsf {T} (\mathbf {M}), u ^ {*}, v ^ {*} \quad \text {in} \quad \mathsf {T} (\mathbf {M} ^ {*}) \\ \langle t \llcorner u ^ {*}, v ^ {*} \rangle = \langle t, \theta_ {\mathsf {S}} (u ^ {*} v ^ {*}) \rangle & \text {for} \quad t \in \mathsf {S} (\mathbf {M}), u ^ {*}, v ^ {*} \quad \text {in} \quad \mathsf {S} (\mathbf {M} ^ {*}) \\ \langle v ^ {*}, u ^ {*} \lrcorner t \rangle = \langle \theta_ {\wedge} (u ^ {*} \wedge v ^ {*}), t \rangle & \text {for} \quad t \in \bigwedge (\mathbf {M}), u ^ {*}, v ^ {*} \quad \text {in} \quad \bigwedge (\mathbf {M} ^ {*}). \end{array}
\]

我们把显式求出左内乘的类似公式的任务留给读者，这次使用公式 (61)、(62) 和 (63)。

上述内容可以在把 M 换成其对偶 \(M^{*}\) 的情形下应用；此时必须把 \(M^{*}\) 换成双重对偶 \(M^{**}\) ，例如， \(T(M^{*})\) 于是具有代数 \(T(M^{**})\) 上的一个右模结构。但典范映射 \(c_{M}: M \to M^{**}\) 定义了代数同态 \(T(c_{M}): T(M) \to T(M^{**})\) ，借助它， \(T(M^{*})\) 具有一个右 \(T(M)\) -模结构。类似地， \(S(M^{*})\) （相应地 \(\bigwedge(M^{*})\) ）具有一个右 \(S(M)\) -模（相应地，左 \(\bigwedge(M)\) -模）结构。给出这些模的外合成律的显式公式可通过交换 M 与 \(M^{*}\) 的角色从上文直接得出。注意，对所有 \(x \in M\) , \(i(x)\) 总是分次代数 \(S(M^{*})\) （相应地 \(\bigwedge(M^{*})\) ）的导子（相应地，平方为零的反导子）。

命题 11. 典范同态 \(\theta_{\mathsf{T}}: \mathsf{T}(\mathbf{M}^{*}) \to \mathsf{T}(\mathbf{M})^{*\mathsf{gr}}\) （相应地 \(\theta_{\mathsf{S}}: \mathsf{S}(\mathbf{M}) \to \mathsf{S}(\mathbf{M})^{*\mathsf{gr}}\) ，相应地 \(\theta_{\wedge}: \bigwedge (\mathbf{M}^{*}) \to \bigwedge (\mathbf{M})^{*\mathsf{gr}}\) ）是一个右 \(\mathsf{T}(\mathbf{M})\) -模（相应地，右 \(\mathsf{S}(\mathbf{M})\) -模，相应地，左 \(\bigwedge (\mathbf{M})\) -模）同态。

我们首先证明，对 \(z^{*} \in \mathsf{T}(\mathbf{M}^{*})\) 和 \(t \in \mathsf{T}(\mathbf{M})\) ，

\[
\theta_ {T} (z ^ {*} \llcorner t) = \theta_ {T} (z ^ {*}) \llcorner t.\tag{69}
\]

由于 M 是代数 \(\mathsf{T}(\mathbf{M})\) 的生成系，我们只需在 \(t = x \in M\) 时证明 (69)；此外，我们可以把注意力限制在 \(z^{*} = x_{1}^{*} \otimes x_{2}^{*} \otimes \cdots \otimes x_{p}^{*}\) （其中 \(x_{j}^{*} \in M^{*}\) ）的情形；于是，由 (66)（交换 M 与 \(M^{*}\) 的角色）， \(z^{*} \llcorner x = \langle x, x_{1}^{*} \rangle x_{2}^{*} \otimes \cdots \otimes x_{p}^{*}\) 。因此，对 M 中所有 \(y_{2}, \ldots, y_{p}\) ，

\[
\begin{array}{r l} \langle \theta_ {T} (z ^ {*} \llcorner x), y _ {2} \otimes y _ {3} \otimes \dots \otimes y _ {p} \rangle & = \langle x, x _ {1} ^ {*} \rangle \prod_ {j = 2} ^ {p} \langle y _ {j}, x _ {j} ^ {*} \rangle \\ & = \langle \theta_ {T} (z ^ {*}), x \otimes y _ {2} \otimes \dots \otimes y _ {p} \rangle \\ & = \langle \theta_ {T} (z ^ {*}) \llcorner x, y _ {2} \otimes \dots \otimes y _ {p} \rangle \end{array}
\]

由此即得 (69)。

其次我们证明对 \(z^{*} \in \mathbf{S}(\mathbf{M}^{*})\) 和 \(t \in \mathbf{S}(\mathbf{M})\) ,

\[
\theta_ {S} (z ^ {*} \llcorner t) = \theta_ {S} (z ^ {*}) \llcorner t.\tag{70}
\]

如上所述，我们可以限于 \(t = x \in \mathbf{M}\) 的情形。但进一步地，这里 \(i(x)\) 是 \(\mathbf{S}(\mathbf{M}^*)\) 的导子，也是 \(\mathbf{S}(\mathbf{M})^{*\text{gr}}\) 的导子。因此（§ 10，第 7 号，命题 9 的推论），只需对 \(z^* = x^* \in \mathbf{M}^*\) 验证 (70)，因为 \(\mathbf{M}^*\) 是 \(\mathbf{S}(\mathbf{M}^*)\) 的生成系；但这是平凡的，此时两边都等于 \(\langle x^*, x \rangle\) 。类似的论证可证明关系式

\[
\theta_ {\wedge} (t \lrcorner z ^ {*}) = t \lrcorner \theta_ {\wedge} (z ^ {*})\tag{71}
\]

对 \(z^{*} \in \bigwedge(\mathbf{M}^{*})\) 和 \(t \in \bigwedge(\mathbf{M})\) ：此时注意到，对 \(x \in \mathbf{M}\) , \(i(x)\) 在 \(\bigwedge(\mathbf{M}^{*})\) 中以及 \(\bigwedge(\mathbf{M})^{*\text{er}}\) 中都是反导子，并使用 § 10，第 7 号，命题 9 的推论。对左内乘有类似的结果。

